\section{Scope, provenance, and principal conclusions}
\label{agrc:scope}

The repository's polylogarithm programme already includes a substantial
calculus of Hurwitz jets, harmonic Dirichlet series, Gamma quotients,
Stieltjes constants, and regularized correlations. Its current reading
artifact distinguishes accepted scientific material from queued incoming
reports \cite{repo-manuscript}. That distinction is retained here.

The most specific problem addressed is the request in the
\emph{Stieltjes--harmonic identities} research agenda for a direct finite
Bernoulli-polynomial proof of the Gauss residue identity
\cite{repo-agenda,repo-gauss}. The identity was already proved there
analytically; the open task was an independent algebraic explanation.
Theorem~\ref{agrc:residue} supplies that explanation. More importantly,
Theorem~\ref{agrc:tail} evaluates every coefficient after the residue,
not only the residue itself. This is the main algebraic extension.

The harmonic part starts from a different preserved report, whose
all-depth, outer-shifted Laurent theorem already determines the principal
and constant terms at nonpositive integers \cite{repo-laurent}. We use
that theorem as established background and give a short independent
local argument where needed. The extension developed here is the
\emph{moving} reflection and centering: the shift varies with the
harmonic-order parameter before coefficients are extracted. This
produces common holomorphic zeros of all depths, an elementary
uniqueness theorem, a Gamma generating function for convergent sums,
and an unequal-frequency multivariate version.

The generating Barnes--Hurwitz family itself is classical. Young's
height-one generating theorem \cite{Young2024} is a direct antecedent;
related Laurent and harmonic-zeta work includes
\cite{Young2023,Can2024}. Standard Bernoulli and Hurwitz conventions are
those in \cite{DLMF24,DLMF25}. We do not infer global originality from
failure to locate an identical formula. The new statements are presented
as proved continuations relative to the inspected textual sources, with
an explicit overlap map in the accompanying audit.

\subsection{What is proved, and what is not}

The article establishes the following interconnected conclusions.
\begin{enumerate}[label=\arabic*.]
\item An elementary residue proof and a finite formula for the full
resonant Gauss tail, including shift transport and mixed parameter jets.
\item A moving-reflection theorem, common all-depth trivial zeros,
and uniqueness of the centered affine shift.
\item An every-depth zeta-cumulant evaluation of convergent harmonic sums;
a compatible centered Stieltjes normalization; and all-order shifted
polygamma--logarithm identities with ordinary antiderivatives.
\item A reflection theorem for any finite number of positive, unequal
frequencies and independent harmonic-depth parameters.
\end{enumerate}
No claim is made that individual colored polylogarithms reduce to depth
one, that all higher harmonic Stieltjes coefficients have finite Gamma
reductions, or that S6 or S8 vanishes. The incoming directory was inspected
as a queue; its binary archives were not all unpacked and audited in this
run. This is not a completeness claim about that queue.

\section{Conventions and the harmonic Barnes family}
\label{agrc:preliminaries}

For $n,r\geq0$, put
\begin{equation}\label{agrc:elementary}
 E_{n,r}=[u^r]\prod_{j=1}^{n}\left(1+\frac uj\right),
 \qquad E_{n,0}=1.
\end{equation}
An empty product is one. Thus
\begin{align}
 E_{n,1}&=H_n,\notag\\
 E_{n,2}&=\frac{H_n^2-H_n^{(2)}}2,\notag\\
 E_{n,3}&=\frac{H_n^3-3H_nH_n^{(2)}+2H_n^{(3)}}6,
 \qquad H_n^{(j)}=\sum_{k=1}^n k^{-j}.\label{agrc:elementary-examples}
\end{align}
With $\Re a>0$, define
\begin{equation}\label{agrc:BZ}
 \Bcal(s,u;a)=\sum_{n=0}^{\infty}
       \frac{(1+u)_n}{n!(n+a)^s},
 \qquad Z_r(s,a)=[u^r]\Bcal(s,u;a).
\end{equation}
The series for $Z_r$ converges normally for $\Re s>1$; the bivariate
series converges normally on compact subsets of, for example,
$|u|<\delta$, $\Re s>1+\delta$, $\Re a>0$.
All powers of $n+a$ use the principal logarithm. At $a=1$,
$Z_r(s,1)=\zeta(s,\{1\}^r)$ in the convention with the largest index first.

\begin{lemma}[Mellin representation and uniform subtraction]
\label{agrc:mellin}
In the common convergence domain,
\begin{equation}\label{agrc:mellin-eq}
 \Gamma(s)\Bcal(s,u;a)
 =\int_0^\infty t^{s-1}\e^{-at}(1-\e^{-t})^{-1-u}\dd t.
\end{equation}
Define the analytic heat coefficients at $t=0$ by
\begin{equation}\label{agrc:heat}
 \phi(t,u,a)=\e^{-at}\left(\frac{t}{1-\e^{-t}}\right)^{1+u}
       =\sum_{j\geq0}c_j(u,a)t^j.
\end{equation}
For every $K\geq1$,
\begin{equation}\label{agrc:subtraction}
 \Gamma(s)\Bcal(s,u;a)
  =\sum_{j=0}^{K-1}\frac{c_j(u,a)}{s-1-u+j}+J_K(s,u;a),
\end{equation}
where $J_K$ is jointly holomorphic for $\Re(s-u)>1-K$, locally in
$\Re a>0$. Increasing $K$ continues the same meromorphic family.
\end{lemma}
\begin{proof}
The binomial series is
$\sum_{n\geq0}(1+u)_n\e^{-nt}/n!=(1-\e^{-t})^{-1-u}$.
Termwise integration is justified by absolute convergence and the
Laplace integral for $(n+a)^{-s}$. Split the integral at $t=1$.
On $(0,1)$, subtract the first $K$ terms of \eqref{agrc:heat}; the
remainder is $O(t^K)$ uniformly on compact parameter sets.
The remaining small-$t$ integrand is bounded by a constant times
$t^{\Re(s-u)-2+K}$. Parameter derivatives add only fixed powers of
$\log t$ and preserve integrability on smaller compact subsets.
At infinity the positive real part of $a$ gives exponential decay,
also after differentiation. Integrating the subtracted monomials gives
\eqref{agrc:subtraction}. The formulas agree in the original half-plane,
so their continuations are compatible.
\end{proof}

The heat coefficients are rational polynomials in $u,a$. A convenient
exact recursion follows from
\begin{equation}\label{agrc:heat-log}
 \log\phi(t,u,a)=\left(\frac{1+u}{2}-a\right)t
 -(1+u)\sum_{k\geq1}\frac{B_{2k}}{2k(2k)!}t^{2k}.
\end{equation}
This is an analytic germ, not an asymptotic series. The Bernoulli
numbers use $B_1=-1/2$.

We will also use the standard Laurent convention
\begin{equation}\label{agrc:Stieltjes}
 \zeta(1+\varepsilon,a)=\varepsilon^{-1}
 +\sum_{j\geq0}\frac{(-1)^j\gamma_j(a)}{j!}\varepsilon^j,
 \qquad \gamma_0(a)=-\psi(a).
\end{equation}
A finite part always denotes a Laurent coefficient of a specified
meromorphic function; it is not an ordinary sum of a divergent series.

\section{A finite algebraic proof of the Gauss residue identity}
\label{agrc:gauss-residue-section}

The coefficients in this section are different from the heat coefficients
$c_j$ above. For indeterminates $a,b,s,c$, define
\begin{align}
 D_j(a,b,s;c)=\frac{(-1)^{j+1}}{j(j+1)}
 \bigl(&B_{j+1}(a-c)+B_{j+1}(b-c)\notag\\
       &-B_{j+1}(a+b+s-c)-B_{j+1}(1-c)\bigr),\label{agrc:D}\\
 \sum_{m\geq0}A_m(a,b,s;c)t^m
 &=\exp\left(\sum_{j\geq1}D_j(a,b,s;c)t^j\right).
 \label{agrc:A}
\end{align}
The last identity is formal. In particular,
\begin{equation}\label{agrc:A-rec}
 A_0=1,\qquad mA_m=\sum_{j=1}^m jD_j A_{m-j}.
\end{equation}
For every fixed $m$, this is a finite polynomial definition over $\Q$.

\begin{lemma}[Exact comparison-shift transport]
\label{agrc:transport}
For $m\geq1$,
\begin{equation}\label{agrc:transport-derivative}
 \partial_c A_m(a,b,s;c)=(s+m)A_{m-1}(a,b,s;c).
\end{equation}
Consequently, for an arbitrary shift $h$,
\begin{equation}\label{agrc:transport-finite}
 A_m(a,b,s;c+h)=\sum_{j=0}^{m}
        \binom{s+m}{j}h^j A_{m-j}(a,b,s;c).
\end{equation}
The binomial coefficient here is its polynomial meaning.
\end{lemma}
\begin{proof}
The Bernoulli derivative formula gives
$\partial_cD_1=s+1$ and
$\partial_cD_j=(j-1)D_{j-1}$ for $j\geq2$.
Writing the formal series in \eqref{agrc:A} as $G(t)$ gives
$\partial_cG=(s+1)tG+t^2\partial_tG$. Coefficient comparison proves
\eqref{agrc:transport-derivative}. Iterating that derivative and applying
the finite Taylor formula proves \eqref{agrc:transport-finite}.
\end{proof}

At $s=-r$, the coefficient $A_r$ is already seen to be independent of $c$.
The following argument determines it without Gamma functions, infinite
sums, or meromorphic residues.

\begin{theorem}[Algebraic Gauss residue identity]
\label{agrc:residue}
For every integer $r\geq0$ and all complex $a,b,c$,
\begin{equation}\label{agrc:residue-formula}
 \boxed{\displaystyle
 A_r(a,b,-r;c)=\frac{(-1)^r}{r!}(1-a)_r(1-b)_r.}
\end{equation}
\end{theorem}
\begin{proof}
The case $r=0$ is immediate. Suppose $r\geq1$. Each $D_j$ at $s=-r$
has total degree at most $j+1$ in $a,b$, with coefficients in $\Q[c]$.
Every product contributing to $A_r$ therefore has degree at most $2r$.

Fix an integer $B$ with $1\leq B\leq r$ and put $b=B$.
The finite-difference identity
$B_{j+1}(x+M)-B_{j+1}(x)=(j+1)\sum_{h=0}^{M-1}(x+h)^j$
gives
\begin{equation}\label{agrc:D-root}
 D_j(a,B,-r;c)=\frac{(-1)^{j+1}}j
 \left(\sum_{h=1}^{r-B}(a-c-h)^j+
       \sum_{h=1}^{B-1}(h-c)^j\right).
\end{equation}
Exponentiating the formal logarithms yields the finite product
\begin{equation}\label{agrc:collapse}
 \sum_{m\geq0}A_m(a,B,-r;c)t^m
 =\prod_{h=1}^{r-B}(1+(a-c-h)t)
       \prod_{h=1}^{B-1}(1+(h-c)t).
\end{equation}
Its degree is at most $r-1$. Hence $A_r(a,B,-r;c)=0$.
Symmetry in $a,b$ supplies the roots $a=1,\ldots,r$ as well.
Thus $(1-a)_r(1-b)_r$ divides $A_r$ in $\Q[c,a,b]$.
The degree bound makes the quotient a polynomial in $c$ alone.

Only $D_1^r/r!$ can contribute total degree $2r$: using any $D_j$ with
$j\geq2$ makes the degree strictly smaller. The degree-two homogeneous
part of $D_1(a,b,-r;c)$ in $a,b$ is $-ab$. Therefore the coefficient
of $a^rb^r$ in $A_r$ is $(-1)^r/r!$. The coefficient of $a^rb^r$ in
$(1-a)_r(1-b)_r$ is one. This determines the quotient and proves the
identity, including its independence of $c$.
\end{proof}

\begin{remark}
The residue identity in \cite{repo-gauss} was already an analytic theorem.
The proof above resolves the separate finite-algebraic proof request in
\cite{repo-agenda}. The following section goes beyond that request by
finding every coefficient in the resonant tail.
\end{remark}

\section{The entire resonant tail: Stirling formulas and an Appell family}
\label{agrc:tail-section}

Let $\stirling{\ell}{j}$ be a Stirling number of the second kind, with
$\stirling00=1$. Define the finite polynomial
\begin{equation}\label{agrc:Q-definition}
 Q_{r,k}(a,b;c)=\sum_{\ell=0}^k\binom{k}{\ell}c^{k-\ell}
       \sum_{j=0}^{\ell}(-1)^j\stirling{\ell}{j}
                 \frac{(a)_j(b)_j}{(r+1)_j}.
\end{equation}
Because $r+1$ is a positive integer, there is no singular denominator.

\begin{theorem}[Finite formula for every resonant coefficient]
\label{agrc:tail}
For all integers $r,k\geq0$,
\begin{equation}\label{agrc:tail-formula}
 \boxed{\displaystyle
 A_{r+k}(a,b,-r;c)=
 \frac{(-1)^r}{r!}(1-a)_r(1-b)_r Q_{r,k}(a,b;c).}
\end{equation}
Equivalently, the normalized tail has the convergent local exponential
 generating function
\begin{equation}\label{agrc:tail-hypergeometric}
 \sum_{k\geq0}Q_{r,k}(a,b;c)\frac{t^k}{k!}
       =\e^{ct}\,{}_2F_1(a,b;r+1;1-\e^t).
\end{equation}
It satisfies
\begin{equation}\label{agrc:Appell}
 \partial_c Q_{r,k}=kQ_{r,k-1},\qquad
 Q_{r,k}(c+h)=\sum_{j=0}^k\binom{k}{j}h^jQ_{r,k-j}(c).
\end{equation}
\end{theorem}
\begin{proof}
First take $c=0$. The finite Gamma-ratio asymptotic expansion used in
\cite{repo-gauss}, or its formal Bernoulli derivation, identifies
\begin{equation}\label{agrc:formalW}
 \mathscr W(x)=x^{r-1}\sum_{j\geq0}A_j(a,b,-r;0)x^{-j}
\end{equation}
as the formal Laurent expansion of
$\Gamma(x+a)\Gamma(x+b)/[\Gamma(x+1)\Gamma(x+a+b-r)]$.
For parameters in a right half-plane this follows from the usual finite
log-Gamma expansion; each coefficient identity then extends polynomially
to all $a,b$. The Gamma recurrence consequently gives the formal identity
\begin{equation}\label{agrc:W-recurrence}
 (x+1)(x+a+b-r)\mathscr W(x+1)
       =(x+a)(x+b)\mathscr W(x).
\end{equation}
Only this formal recurrence is used in the rest of the proof.

Discard the polynomial part of $\mathscr W$ and send
$x^{-k-1}$ to $t^k/k!$. Write the resulting formal power series as
\begin{equation}\label{agrc:T-borel}
 T(t)=\sum_{k\geq0}A_{r+k}(a,b,-r;0)\frac{t^k}{k!}.
\end{equation}
Multiplication by $x$ becomes $\partial_t$ after polynomial parts are
discarded; replacing $x$ by $x+1$ becomes multiplication by $\e^{-t}$.
These statements follow directly on each monomial $x^{-k-1}$.
Polynomial terms remain polynomial under both operations, so their
omission does not affect the transformed identity. With $D=\partial_t$,
\eqref{agrc:W-recurrence} becomes
\[
 (D+1)(D+a+b-r)(\e^{-t}T)=(D+a)(D+b)T.
\]
After simplification,
\begin{equation}\label{agrc:T-ODE}
 (1-\e^t)T''+
       (a+b-r-1-(a+b)\e^t)T'-ab\e^tT=0.
\end{equation}
Set $z=1-\e^t$. For $T(t)=F(z)$ this is Gauss's equation
\begin{equation}\label{agrc:Gauss-ODE}
 z(1-z)F''+[r+1-(a+b+1)z]F'-abF=0.
\end{equation}
The unique formal solution with constant term $F(0)$ is
$F(0){}_2F_1(a,b;r+1;z)$: comparison of successive coefficients divides
only by $(j+1)(r+1+j)$, which is nonzero. Theorem~\ref{agrc:residue}
determines $F(0)=A_r$. This proves \eqref{agrc:tail-hypergeometric}
at $c=0$, with the residue prefactor restored.

At $s=-r$ and $m=r+k$, \eqref{agrc:transport-finite} has
$\binom{s+m}{j}=\binom{k}{j}$, which vanishes for $j>k$.
Thus comparison-shift transport multiplies the exponential generating
function of the tail by $\e^{ct}$.
Finally,
$(\e^t-1)^j=j!\sum_{\ell\geq j}\stirling{\ell}{j}t^\ell/\ell!$
converts the hypergeometric series into \eqref{agrc:Q-definition}.
This proves every assertion. Local convergence follows, for example,
in a sufficiently small disc where $|1-\e^t|<1$.
\end{proof}

This is an exponential generating function for the normalized tail,
not a claim of convergence of the original inverse-power expansion
\eqref{agrc:formalW}. The latter remains a formal asymptotic series.
The explicit polynomials begin with
\begin{align}
 Q_{r,0}&=1,\notag\\
 Q_{r,1}&=c-\frac{ab}{r+1},\notag\\
 Q_{r,2}&=c^2-\frac{2ab}{r+1}c
      +\frac{ab(ab+a+b-r-1)}{(r+1)(r+2)}.
 \label{agrc:Q-examples}
\end{align}
For $r=1$ this gives
\[
 A_2(a,b,-1;c)=-(1-a)(1-b)\left(c-\frac{ab}{2}\right).
\]
Unlike a recursive Bernoulli construction, \eqref{agrc:Q-definition}
has a finite triangular formula whose summation ranges are immediately
visible. It also gives the ordinary polynomial primitives
$\int Q_{r,k}(c)\dd c=Q_{r,k+1}(c)/(k+1)+\mathrm{constant}$.

\subsection{Integer-root collapse and nonsingular parameter jets}

Equation~\eqref{agrc:collapse} proves more than the vanishing of $A_r$:
if one upper parameter lies in $\{1,\ldots,r\}$, every $A_m$ with
$m\geq r$ vanishes at $s=-r$. The tail formula exhibits the full common
factor responsible for this collapse.

Write $P_r(a)=(1-a)_r$. At an integer $A\in\{1,\ldots,r\}$, the
following analytic germ avoids logarithmic differentiation through a
zero:
\begin{align}
 P_r(A+x)={}&(-1)^A(A-1)!(r-A)!\,x\notag\\
 &\times\exp\left[-\sum_{j\geq1}
  \frac{H_{r-A}^{(j)}+(-1)^jH_{A-1}^{(j)}}j x^j\right].
 \label{agrc:root-jets}
\end{align}
To prove it, isolate the factor $-x$ in $\prod_{k=1}^r(k-A-x)$;
the other factors are finite products of $1+x/j$ and $1-x/j$.
Thus all mixed jets of \eqref{agrc:tail-formula} at integer roots are
finite harmonic polynomials. For the residue itself,
\begin{equation}\label{agrc:double-root}
 \left.\partial_a\partial_b A_r(a,b,-r;c)\right|_{a=A,b=B}
 =\frac{(-1)^{r+A+B}}{r!}
 (A-1)!(r-A)!(B-1)!(r-B)!.
\end{equation}
No singular polygamma value is involved.

The normal, or $s$-direction, jets are different. For example,
\begin{equation}\label{agrc:normal-transport}
 \left.\partial_c\partial_s^\ell A_r\right|_{s=-r}
 =\left.\ell\,\partial_s^{\ell-1}A_{r-1}\right|_{s=-r}
 \qquad(\ell\geq1).
\end{equation}
Therefore comparison-shift independence of $A_r$ on the resonant
hyperplane does \emph{not} imply independence of its normal derivatives.
The distinction matters in Stieltjes coefficient extraction.

\section{Stieltjes jets and an all-order shifted logarithmic identity}
\label{agrc:Stieltjes-sums}

\subsection{How the tail enters resonant Gauss summation}

For $\Re a,\Re b,\Re c>0$, set
\[
 W_n(a,b;s)=\frac{\Gamma(n+a)\Gamma(n+b)}
                 {\Gamma(n+a+b+s)\Gamma(n+1)}.
\]
The reciprocal Gamma in the denominator is understood as an entire
function. The normally convergent, successively subtracted series
\begin{equation}\label{agrc:Gauss-subtracted}
 \mathcal F^{[M]}(s)=\sum_{n\geq0}
 \left[W_n(a,b;s)-\sum_{j=0}^M A_j(a,b,s;c)(n+c)^{-1-s-j}\right]
\end{equation}
is holomorphic for $\Re s>-M-1$. Its exact value is
\begin{equation}\label{agrc:Gauss-closed}
 \mathcal F^{[M]}(s)=
 \frac{\Gamma(a)\Gamma(b)\Gamma(s)}{\Gamma(a+s)\Gamma(b+s)}
       -\sum_{j=0}^M A_j(a,b,s;c)\zeta(1+s+j,c).
\end{equation}
This is the existing uniform subtraction theorem of \cite{repo-gauss}.
For clarity, its justification is finite-order, locally uniform
log-Gamma asymptotics followed by Gauss summation in $\Re s>0$ and
analytic continuation. The summand in \eqref{agrc:Gauss-subtracted}
is $O(n^{-M-2-\Re s})$ uniformly on compact parameter sets.

Here is the precise coefficient rule, useful for applying the new tail
formula without losing a Stieltjes term. Fix $r\leq M$, and write
\begin{align*}
 A_j(a,b,-r+\varepsilon;c)&=\sum_{\ell\geq0}a_{j,\ell}\varepsilon^\ell,\\
 R(\varepsilon)&=\varepsilon
 \frac{\Gamma(a)\Gamma(b)\Gamma(-r+\varepsilon)}
      {\Gamma(a-r+\varepsilon)\Gamma(b-r+\varepsilon)}
       =\sum_{\ell\geq0}R_\ell\varepsilon^\ell.
\end{align*}
The germ $R$ is analytic, including at zeros of either reciprocal Gamma.
For $k\geq0$, direct Laurent multiplication gives
\begin{align}
 [\varepsilon^k]\mathcal F^{[M]}(-r+\varepsilon)
 ={}&R_{k+1}-a_{r,k+1}
 -\sum_{\ell=0}^k a_{r,\ell}
           \frac{(-1)^{k-\ell}\gamma_{k-\ell}(c)}{(k-\ell)!}\notag\\
 &-\sum_{\substack{0\leq j\leq M\\j\ne r}}\sum_{\ell=0}^k
   a_{j,\ell}\frac{\zeta^{(k-\ell)}(1-r+j,c)}{(k-\ell)!}.
 \label{agrc:Gauss-jets}
\end{align}
The superscript on $\zeta$ denotes spectral differentiation.
Theorem~\ref{agrc:tail} now gives every $a_{r+j,0}$ by a finite Stirling
sum. The coefficients $a_{r+j,\ell}$ for $\ell>0$ still require normal
parameter differentiation; replacing them by tangential derivatives of
the resonant formula would be incorrect.

\subsection{A shifted Beta--Stieltjes generating theorem}

The following especially simple specialization gives ordinary convergent
identities at every logarithmic order. Define
\begin{equation}\label{agrc:Pell}
 P_\ell(x)=[\varepsilon^\ell]\frac{\Gamma(x)}{\Gamma(x+\varepsilon)}
 =\frac1{\ell!}Y_\ell\bigl(-\psi(x),-\psi'(x),\ldots,
                                      -\psi^{(\ell-1)}(x)\bigr),
\end{equation}
where $Y_\ell$ is the complete exponential Bell polynomial and $P_0=1$.
Equivalently, the first equality defines the polynomials without Bell
notation. Set $b_{-1}(x)=1$, $b_j(x)=(-1)^j\gamma_j(x)/j!$ for $j\geq0$,
and $b_j(x)=0$ for $j<-1$.

\begin{theorem}[All shifted logarithmic polygamma sums]
\label{agrc:beta-stieltjes}
For $\Re x>0$, the series
\begin{align}
 \Dcal_x(\varepsilon)=\sum_{n\geq0}\bigg[
 &\frac{\Gamma(n+x)}{\Gamma(n+x+\varepsilon)}
 -(n+x)^{-\varepsilon}\notag\\
 &-\frac{\varepsilon(1-\varepsilon)}2(n+x)^{-1-\varepsilon}\bigg]
 \label{agrc:Dx}
\end{align}
is normally convergent for $\Re\varepsilon>-1$. It equals the
holomorphic completion of
\begin{equation}\label{agrc:Dx-closed}
 \boxed{\displaystyle
 \Dcal_x(\varepsilon)=
 \frac{\Gamma(x)}{(\varepsilon-1)\Gamma(x+\varepsilon-1)}
 -\zeta(\varepsilon,x)
 -\frac{\varepsilon(1-\varepsilon)}2\zeta(1+\varepsilon,x).}
\end{equation}
For every integer $\ell\geq1$, its coefficient identity is
\begin{align}
 &\sum_{n\geq0}\bigg[P_\ell(n+x)
 -\frac{(-\log(n+x))^\ell}{\ell!}\notag\\
 &\qquad-\frac1{2(n+x)}\left\{
       \frac{(-\log(n+x))^{\ell-1}}{(\ell-1)!}
       -\mathbf1_{\ell\geq2}
          \frac{(-\log(n+x))^{\ell-2}}{(\ell-2)!}\right\}\bigg]\notag\\
 &=P_\ell(x)-x\sum_{j=0}^{\ell}P_j(x)
    -\frac{\zeta^{(\ell)}(0,x)}{\ell!}
    -\frac12 b_{\ell-1}(x)+\frac12 b_{\ell-2}(x).
 \label{agrc:all-log-sums}
\end{align}
The grouped series in this formula is absolutely convergent.
\end{theorem}
\begin{proof}
The uniform two-term Gamma-ratio expansion is
\[
 \frac{\Gamma(y)}{\Gamma(y+\varepsilon)}
 =y^{-\varepsilon}\left(1+
         \frac{\varepsilon(1-\varepsilon)}{2y}+O(y^{-2})\right),
 \qquad y=n+x\longrightarrow\infty.
\]
Uniformity on compact parameter sets proves normal convergence of
\eqref{agrc:Dx}; all fixed parameter derivatives remain normally
convergent. In $\Re\varepsilon>1$, the beta integral gives
\begin{align*}
 \sum_{n\geq0}\frac{\Gamma(n+x)}{\Gamma(n+x+\varepsilon)}
 &=\frac1{\Gamma(\varepsilon)}\int_0^1
        t^{x-1}(1-t)^{\varepsilon-1}\sum_{n\geq0}t^n\dd t\\
 &=\frac{\Gamma(x)}{(\varepsilon-1)\Gamma(x+\varepsilon-1)}.
\end{align*}
The two comparison series sum to the displayed Hurwitz terms. Analytic
continuation proves \eqref{agrc:Dx-closed} on the larger half-plane,
including the removable singularities at $\varepsilon=0,1$.

For coefficient extraction, the Gamma recurrence rewrites the first
term of \eqref{agrc:Dx-closed} as
\[
 \left(1-\frac{x}{1-\varepsilon}\right)
          \frac{\Gamma(x)}{\Gamma(x+\varepsilon)}.
\]
Its coefficient of $\varepsilon^\ell$ is
$P_\ell(x)-x\sum_{j=0}^{\ell}P_j(x)$. Multiplication by
$\varepsilon(1-\varepsilon)/2$ in the Hurwitz Laurent expansion gives
$(b_{\ell-1}-b_{\ell-2})/2$. The remaining terms are ordinary Taylor
coefficients. On the series side, every grouped coefficient is
$O(n^{-2}(1+\log^\ell n))$ on compact $x$-sets; hence extraction commutes
with summation and proves \eqref{agrc:all-log-sums}.
\end{proof}

These formulas are consequences of a classical beta summation and
precise subtraction, not claims of new transcendental constants.
Their usefulness is an explicit, all-order normalization connecting
nonlinear harmonic polynomials to Stieltjes data and spectral derivatives.

\subsection{Two explicit identities}

Using $\zeta'(0,x)=\log\Gamma(x)-\tfrac12\log(2\pi)$, the first order is
\begin{align}
 &\sum_{n\geq0}\left[\log(n+x)-\psi(n+x)-\frac1{2(n+x)}\right]\notag\\
 &\quad=\left(x-\frac12\right)\psi(x)-\log\Gamma(x)-x
                         +\frac12+\frac12\log(2\pi).
 \label{agrc:first-log}
\end{align}
The second order is
\begin{align}
 &\sum_{n\geq0}\left[\psi(n+x)^2-\psi'(n+x)-\log^2(n+x)
                       +\frac{1+\log(n+x)}{n+x}\right]\notag\\
 &\quad=(1-x)\bigl(\psi(x)^2-\psi'(x)\bigr)
     +(2x-1)\psi(x)-2x+\gamma_1(x)-\zeta''(0,x).
 \label{agrc:second-log}
\end{align}
At $x=1$, these become
\begin{align}
 \sum_{k\geq1}\left[\log k-H_{k-1}+\gamma-\frac1{2k}\right]
 &=\frac{\log(2\pi)-1-\gamma}{2},\label{agrc:first-at-one}\\
 \sum_{k\geq1}\left[(H_{k-1}-\gamma)^2+H_{k-1}^{(2)}-\zeta(2)
            -\log^2 k+\frac{1+\log k}{k}\right]
 &=-2-\gamma+\gamma_1-\zeta''(0).
 \label{agrc:second-at-one}
\end{align}
The grouping is essential. These are not identities obtained by
separately summing their generally divergent components.

\section{Moving reflection, common zeros, and unique centering}
\label{agrc:reflection-section}

The local identity underlying the reflection theorem is
\begin{equation}\label{agrc:heat-reflection}
 c_j(u,1+u-a)=(-1)^jc_j(u,a).
\end{equation}
Indeed,
$\phi(t,u,1+u-a)=\phi(-t,u,a)$ as analytic germs. One can also see it
immediately from \eqref{agrc:heat-log}. The term $u$ in the reflected
shift cannot be dropped.

\begin{theorem}[Moving-reflection zeros at every depth]
\label{agrc:moving-reflection}
Let $0<\Re a<1$ and take $u$ in a sufficiently small neighbourhood of
zero. Define
\[
 \mathcal R_\pm(s,u;a)=\Bcal(s,u;a)\pm\Bcal(s,u;1+u-a).
\]
For every $m\geq0$, $\mathcal R_+$ is jointly holomorphic near
$(s,u)=(-2m,0)$ and
\begin{equation}\label{agrc:Rplus-zero}
 \mathcal R_+(-2m,u;a)=0.
\end{equation}
Similarly, $\mathcal R_-$ is jointly holomorphic near
$(-2m-1,0)$ and vanishes at $s=-2m-1$.
Its depth coefficients are the explicit finite combinations
\begin{equation}\label{agrc:Rcoeff}
 [u^r]\mathcal R_\pm(s,u;a)=
 Z_r(s,a)\pm\sum_{k=0}^r\frac{(-1)^k(s)_k}{k!}
                                  Z_{r-k}(s+k,1-a).
\end{equation}
Every coefficient inherits the stated holomorphic zero.
\end{theorem}
\begin{proof}
Apply Lemma~\ref{agrc:mellin} to both factors with the reflected shift
inserted before coefficient extraction. By \eqref{agrc:heat-reflection},
the plus combination has only even heat indices $j$, while the minus
combination has only odd indices. Near $s=-2m,u=0$, the only potentially
vanishing denominator in \eqref{agrc:subtraction} has $j=2m+1$; its
coefficient in the plus combination is identically zero in $u$.
Thus $\Gamma(s)\mathcal R_+$ is jointly holomorphic there. The reciprocal
Gamma has a zero at $s=-2m$, proving the holomorphy and the zero of
$\mathcal R_+$ itself. The minus case uses the excluded even index
$j=2m+2$ in exactly the same argument.

For the coefficient formula, Taylor-expand the shift:
$\partial_a^k\Bcal(s,u;a)=(-1)^k(s)_k\Bcal(s+k,u;a)$
in the ordinary convergence half-plane. Taking coefficients gives
\eqref{agrc:Rcoeff}; meromorphic continuation extends the identity.
Joint holomorphy at the indicated points permits coefficient extraction
there and transfers the zero to every depth.
\end{proof}

\subsection{The centered family}

Put
\begin{equation}\label{agrc:center-family}
 \Ccal(s,u)=\Bcal\left(s,u;\frac{1+u}{2}\right),
 \qquad C_r(s)=[u^r]\Ccal(s,u).
\end{equation}
Its Mellin kernel is particularly simple:
\begin{equation}\label{agrc:center-Mellin}
 \Gamma(s)\Ccal(s,u)=\int_0^\infty
               t^{s-1}\bigl(2\sinh(t/2)\bigr)^{-1-u}\dd t.
\end{equation}
The heat factor
$(t/(2\sinh(t/2)))^{1+u}$ is even.

\begin{corollary}[Centered all-depth trivial zeros]
\label{agrc:center-zeros}
For all integers $r,m\geq0$, the meromorphic combination
\begin{equation}\label{agrc:center-coeff}
 \boxed{\displaystyle
 C_r(s)=\sum_{k=0}^r\frac{(-1)^k(s)_k}{2^k k!}
                               Z_{r-k}(s+k,1/2)}
\end{equation}
is holomorphic at $s=-2m$ and satisfies
\begin{equation}\label{agrc:center-zero}
 \boxed{C_r(-2m)=0.}
\end{equation}
\end{corollary}
\begin{proof}
The even heat factor gives the same missing-pole argument as
Theorem~\ref{agrc:moving-reflection}. Taylor expansion of the shift by
$u/2$ gives \eqref{agrc:center-coeff}.
\end{proof}

The identity is about the continued \emph{combination}. Individual terms
on its right can have poles, and they must not be evaluated separately
at a singular point. At depth zero it reduces to the classical identity
$C_0(s)=\zeta(s,1/2)=(2^s-1)\zeta(s)$ and its even zeros.
At depth one,
\begin{equation}\label{agrc:C1}
 C_1(s)=Z_1(s,1/2)-\frac{s}{2}\zeta(s+1,1/2).
\end{equation}
Consequently $Z_1(0,1/2)=1/2$, and for $m\geq1$,
\begin{equation}\label{agrc:Z1-negative-even}
 Z_1(-2m,1/2)=-m\zeta(1-2m,1/2)=\frac{B_{2m}(1/2)}2.
\end{equation}
Here the individual $Z_1$ values are holomorphic.

At depth two and $m\geq1$, expansion of \eqref{agrc:center-coeff} gives
\begin{align}
 \Res_{s=-2m}Z_2(s,1/2)&=-m\zeta(1-2m,1/2),\label{agrc:Z2-residue}\\
 \FP_{s=-2m}Z_2(s,1/2)
 &=-m\FP_{s=1-2m}Z_1(s,1/2)
                         +\tfrac12\zeta(1-2m,1/2).
 \label{agrc:Z2-fp}
\end{align}
The last $\tfrac12\zeta$ term comes from the first-order variation of
$-s/2$ multiplying a pole. The remaining Hurwitz term is zero because
$\zeta(2-2m,1/2)=0$. For example,
\begin{equation}\label{agrc:Z2-example}
 Z_2(-2+\varepsilon,1/2)
      =-\frac1{24\varepsilon}+\frac{1-2\gamma}{48}+O(\varepsilon).
\end{equation}
The local coefficient rule of \cite{repo-laurent}, reproduced briefly in
Section~\ref{agrc:audit-section}, supplies these individual Laurent data.

\subsection{Why the center is forced}

\begin{theorem}[Unique centering; two-depth rigidity for affine shifts]
\label{agrc:unique-center}
Let $a(u)$ be analytic near zero with $\Re a(0)>0$.
Every coefficient of $\Bcal(s,u;a(u))$ in $u$ is holomorphic and zero at
$s=0$ if and only if
\begin{equation}\label{agrc:unique-shift}
 a(u)=\frac{1+u}{2}
\end{equation}
as an analytic germ. In the affine class $a(u)=a_0+\lambda u$, it is
enough to require that the depth-zero and depth-one coefficients be
holomorphic and zero at $s=0$.
\end{theorem}
\begin{proof}
The resonant heat coefficient at $s=0$ is
$q(u)=c_1(u,a(u))=(1+u)/2-a(u)$. Lemma~\ref{agrc:mellin} gives
\[
 \Bcal(s,u;a(u))=\frac1{\Gamma(s)}\left(\frac{q(u)}{s-u}+Q(s,u)\right),
\]
with $Q$ jointly holomorphic near $(0,0)$. Since
$1/\Gamma(s)=s+O(s^2)$, the $Q$ term is $O(s)$ after any fixed depth
extraction. Suppose the first nonzero coefficient of $q$ is $q_j u^j$.
The depth-$j$ coefficient of the displayed singular term then has
constant term $q_j$ at $s=0$, contradicting the hypothesis. Thus $q=0$.
The converse is Corollary~\ref{agrc:center-zeros}.

For an affine shift,
$q(u)=1/2-a_0+(1/2-\lambda)u$. The depth-zero condition forces
$a_0=1/2$; the depth-one condition then forces $\lambda=1/2$.
All higher even zeros follow automatically.
\end{proof}

This is a finite rigidity phenomenon: in the affine class, two local
conditions select a family with infinitely many zero identities at
all harmonic depths.

\section{An all-depth Gamma evaluation of ordinary harmonic sums}
\label{agrc:ordinary-sums}

Let $L=\log2$, and introduce the centered Gamma quotient
\begin{equation}\label{agrc:Fc}
 F_c(u)=\frac{\Gamma((1+u)/2)\Gamma(1-u)}{\Gamma((1-u)/2)}.
\end{equation}
The subscript denotes ``centered'', not the comparison shift of the
Gauss coefficients.

\begin{theorem}[Centered weight-two Gamma generating function]
\label{agrc:Gamma-center}
For $|u|<1$, with the removable value at $u=0$ understood,
\begin{equation}\label{agrc:Gamma-center-eq}
 \boxed{\displaystyle
 \Ccal(2,u)=\frac{\pi\tan(\pi u/2)}{u}\,F_c(u).}
\end{equation}
Equivalently,
\begin{equation}\label{agrc:cumulant-gen}
 \Ccal(2,u)=\frac{\pi^2}{2}
       \exp\left(\sum_{j\geq1}\frac{\kappa_j}{j}u^j\right),
\end{equation}
where
\begin{equation}\label{agrc:kappa}
 \kappa_1=-2L,\qquad
 \kappa_j=
 \begin{cases}
 (3-2^{2-j})\zeta(j),&j\geq2\text{ even},\\
 (2^{1-j}-1)\zeta(j),&j\geq3\text{ odd}.
 \end{cases}
\end{equation}
If $h_0=1$ and $rh_r=\sum_{j=1}^r\kappa_jh_{r-j}$, then
$C_r(2)=\pi^2h_r/2$ for every $r\geq0$.
\end{theorem}
\begin{proof}
For $\Re u<0$, beta integration at $s=1$ gives
$\Bcal(1,u;a)=B(a,-u)$. Differentiating in $a$ gives
\[
 \Bcal(2,u;a)=-B(a,-u)\bigl(\psi(a)-\psi(a-u)\bigr).
\]
At $a=(1+u)/2$, one has $a-u=1-a$, and the digamma reflection formula
implies
$\psi(a)-\psi(1-a)=-\pi\cot(\pi a)=\pi\tan(\pi u/2)$.
Also $B(a,-u)=-F_c(u)/u$. This proves
\eqref{agrc:Gamma-center-eq} initially in a left half-disc. Both sides
are analytic in $|u|<1$ after removing the value at zero: on the series
side the summands are $O(n^{-2+\Re u})$, locally uniformly, and the
moving shift retains positive real part. Analytic continuation proves
the stated domain.

Taylor expansion of the three log-Gamma factors gives
\begin{equation}\label{agrc:Fc-log}
 \log F_c(u)=-2Lu+
 \sum_{\substack{j\geq2\\j\text{ even}}}\frac{\zeta(j)}j u^j
 +\sum_{\substack{j\geq3\\j\text{ odd}}}
          \frac{(2^{1-j}-1)\zeta(j)}j u^j.
\end{equation}
Here the odd Taylor coefficients at $1/2$ use
$\psi^{(j-1)}(1/2)=(-1)^j(j-1)!(2^j-1)\zeta(j)$ for $j\geq2$;
the linear term uses $\psi(1/2)=-\gamma-2L$.
The sine product, or its Taylor series, yields
\[
 \log\frac{\tan(\pi u/2)}{\pi u/2}
 =\sum_{k\geq1}\frac{(1-2^{1-2k})\zeta(2k)}k u^{2k}.
\]
Adding this to \eqref{agrc:Fc-log} gives \eqref{agrc:kappa}.
Coefficient extraction from the exponential gives the finite recursion.
\end{proof}

\subsection{Explicit generalized harmonic identities}

The coefficient of $(1+u/(2n+1))^{-2}$ gives the following ordinary
series interpretation of every $h_r$:
\begin{equation}\label{agrc:harmonic-all-depth}
 \boxed{\displaystyle
 \sum_{n\geq0}\frac1{(2n+1)^2}
    \sum_{j=0}^r\frac{(-1)^j(j+1)E_{n,r-j}}{(2n+1)^j}
       =\frac{\pi^2}{8}h_r.}
\end{equation}
The $n=0$ terms are included. The bound
$E_{n,r}=O_r((1+\log(n+1))^r)$ proves absolute convergence for each
fixed depth. This is a family of identities, not an assertion that an
arbitrary nonlinear harmonic numerator has the same reduction.

The first coefficients are
\begin{align}
 h_0&=1,&h_1&=-2L,\notag\\
 h_2&=2L^2+\zeta(2),&
 h_3&=-\frac43L^3-2L\zeta(2)-\frac14\zeta(3),\notag\\
 h_4&=\frac23L^4+2L^2\zeta(2)+\frac12L\zeta(3)
                       +\frac12\zeta(2)^2+\frac{11}{16}\zeta(4).
 \label{agrc:h-examples}
\end{align}
For example, writing $q_n=2n+1$,
\begin{align}
 \sum_{n\geq0}\frac{H_n-2/q_n}{q_n^2}
   &=-\frac{\pi^2}{4}L,\label{agrc:h-linear}\\
 \sum_{n\geq0}\frac{H_n^2-H_n^{(2)}-4H_n/q_n+6/q_n^2}{q_n^2}
   &=\frac{\pi^2}{4}\bigl(2L^2+\zeta(2)\bigr),
 \label{agrc:h-quadratic}\\
 \sum_{n\geq0}\frac{E_{n,3}-2E_{n,2}/q_n+3H_n/q_n^2-4/q_n^3}{q_n^2}
   &=\frac{\pi^2}{8}\left(-\frac43L^3-2L\zeta(2)-\frac14\zeta(3)\right).
 \label{agrc:h-cubic}
\end{align}
The complete recurrence, rather than only these small cases, is the
content of \eqref{agrc:harmonic-all-depth}. The data file
\texttt{exact\_tables.json} supplies coefficients through depth eight.

There is a second, independent integral interpretation:
\begin{equation}\label{agrc:log-sinh-moments}
 \frac1{r!}\int_0^\infty
 \frac{t\,[-\log(2\sinh(t/2))]^r}{2\sinh(t/2)}\dd t
     =\frac{\pi^2}{2}h_r.
\end{equation}
At zero the integrand is $O(1+|\log t|^r)$; at infinity it decays
exponentially times a polynomial. Thus this is an ordinary integral.
Differentiation of \eqref{agrc:center-Mellin} at $s=2$ is justified
by these majorants. These moments can change sign; no positivity claim
is being used.

\subsection{All positive outer exponents}

Let
\[
 F(a,u)=\frac{\Gamma(a)\Gamma(1-u)}{\Gamma(a-u)},
 \qquad \Ccal(s,u;x)=\Bcal\left(s,u;\frac{1+u}{2}+x\right).
\]
For every integer $p\geq2$ and $|\Re x|<1/2$, one obtains
\begin{equation}\label{agrc:positive-all-p}
 [u^r]\Ccal(p,u;x)=
 \frac{(-1)^p}{(p-1)!}[u^{r+1}]
  \left.\partial_a^{p-1}F(a,u)\right|_{a=(1+u)/2+x}.
\end{equation}
Indeed,
$\Bcal(p,u;a)=(-1)^{p-1}\partial_a^{p-1}B(a,-u)/(p-1)!$;
substitute $B(a,-u)=-F(a,u)/u$ and extract the coefficient.
The beta identity first holds for $\Re u<0$ and then extends through
zero after differentiation in $a$. Uniform convergence of the differentiated
series for $|u|$ small supplies the same conclusion directly.

For $x=0$ every coefficient in \eqref{agrc:positive-all-p} is a finite
polynomial in $L$ and zeta values at positive integers. Formula
\eqref{agrc:Gamma-center-eq} is the particularly simple reflection
collapse at $p=2$; the other exponents retain finite polygamma/Bell
expressions. This general positive-integer beta mechanism is classical
and overlaps the fixed-shift theory in \cite{repo-laurent}.

\section{Centered Stieltjes constants and polylogarithmic coordinates}
\label{agrc:center-Stieltjes-section}

Write
\begin{equation}\label{agrc:Gf}
 G(u)=\frac1{\Gamma(1+u)}=\sum_{j\geq0}g_ju^j,
 \qquad F_c(u)=\sum_{j\geq0}f_ju^j.
\end{equation}
Thus $g_0=1$, $g_1=\gamma$, and
$g_2=(\gamma^2-\zeta(2))/2$.

\begin{theorem}[Centered principal part and finite part]
\label{agrc:center-Stieltjes}
For every $r\geq0$, as $\varepsilon\to0$,
\begin{equation}\label{agrc:center-Stieltjes-formula}
 \boxed{\displaystyle
 C_r(1+\varepsilon)=
 \sum_{j=0}^r\frac{g_j}{\varepsilon^{r+1-j}}
                +g_{r+1}-f_{r+1}+O(\varepsilon).}
\end{equation}
In particular, the finite part is an explicit polynomial in
$\gamma,L,\zeta(2),\ldots,\zeta(r+1)$.
\end{theorem}
\begin{proof}
The only moving pole near $(s,u)=(1,0)$ has residue
$1/\Gamma(1+u)=G(u)$, because the leading heat coefficient is one.
Consequently
\[
 \Ccal(s,u)=\frac{G(u)}{s-1-u}+Q(s,u)
\]
with $Q$ jointly holomorphic. This follows from
Lemma~\ref{agrc:mellin}, replacing the numerator $1/\Gamma(s)$ by
$1/\Gamma(1+u)$; their difference divided by $s-1-u$ is holomorphic.
For $\Re u<0$, beta integration gives
$\Ccal(1,u)=-F_c(u)/u$, and therefore
$Q(1,u)=(G(u)-F_c(u))/u$. Both sides extend through $u=0$.
Expanding the simple moving pole for $|u|<|\varepsilon|$ and then
continuing the coefficient functions proves
\eqref{agrc:center-Stieltjes-formula}.
\end{proof}

For the first three depths,
\begin{align}
 \FP_{s=1}C_0(s)&=\gamma+2L,\notag\\
 \FP_{s=1}C_1(s)&=\frac{\gamma^2}{2}-\zeta(2)-2L^2,\notag\\
 \FP_{s=1}C_2(s)&=g_3+\frac43L^3+L\zeta(2)+\frac14\zeta(3).
 \label{agrc:center-FP-examples}
\end{align}
These are finite parts of the centered family, not the same constants
as those for a fixed outer shift $1/2$ at positive harmonic depth.

\subsection{Every higher Laurent coefficient: exact transport, not false closure}

Define fixed-shift harmonic Stieltjes coefficients by
\begin{equation}\label{agrc:harmonic-Stieltjes}
 Z_r(1+\varepsilon,a)=
 \sum_{j=0}^r g_j\varepsilon^{-r-1+j}
 +\sum_{\ell\geq0}\frac{(-1)^\ell\gamma_{r,\ell}(a)}{\ell!}
                                             \varepsilon^\ell.
\end{equation}
Define $\widetilde\gamma_{r,\ell}$ in the same way for $C_r$.
At depth zero these are the usual $\gamma_\ell(1/2)$.
For $e_j(k)=[v^j]\prod_{h=1}^k(1+v/h)$, expansion of
\eqref{agrc:center-coeff} gives the exact finite transport law
\begin{align}
 \widetilde\gamma_{r,\ell}
 ={}&\gamma_{r,\ell}(1/2)
 +(-1)^\ell\ell!\sum_{k=1}^r\frac{(-1)^k}{2^k}
       \sum_{j=0}^{\ell}e_j(k)
     \frac{\partial_s^{\ell-j}Z_{r-k}(k+1,1/2)}{(\ell-j)!}.
 \label{agrc:Stieltjes-transport}
\end{align}
Here $e_j(k)=0$ for $j>k$. To prove the formula, observe that
$(1+\varepsilon)_k/k!=\sum_j e_j(k)\varepsilon^j$ and that every
$Z_{r-k}(k+1+\varepsilon,1/2)$ with $k\geq1$ is holomorphic.
The principal part is unchanged, and ordinary Taylor multiplication
proves the identity.

For $\ell=0$, all terms reduce by beta differentiation and reproduce
Theorem~\ref{agrc:center-Stieltjes}. For $\ell>0$, genuine spectral
derivatives remain. Neither centering nor the common zeros justify
discarding them. This is precisely the separation between local finite
parts and higher harmonic Stieltjes data required in the source theory.

The ordinary depth-zero constants satisfy the familiar half-shift law
\begin{equation}\label{agrc:half-Stieltjes}
 \gamma_\ell(1/2)=\gamma_\ell
 +\frac{2(-1)^\ell L^{\ell+1}}{\ell+1}
 +2\sum_{j=0}^{\ell-1}\binom{\ell}{j}
                       (-1)^{\ell+j}L^{\ell-j}\gamma_j.
\end{equation}
This follows directly by expanding
$\zeta(s,1/2)=(2^s-1)\zeta(s)$ at $s=1$.

\subsection{Exact polylogarithmic coordinates at rational shifts}

For $x=p/q$, $1\leq p\leq q$, and $\omega=\e^{2\pi i/q}$,
finite character filtering gives
\begin{equation}\label{agrc:Hurwitz-Li}
 \zeta(s,p/q)=q^{s-1}\sum_{j=0}^{q-1}\omega^{-pj}\Li_s(\omega^j).
\end{equation}
It is first proved by absolutely convergent series in $\Re s>1$ and
then continued meromorphically. In the $j=0$ term, $\Li_s(1)$ means
$\zeta(s)$; the other root-of-unity values use their radial continuation.
At $s=0$ there is no pole, and every spectral derivative has the finite
formula
\begin{align}
 \zeta^{(\ell)}(0,p/q)
  =\frac1q\sum_{j=0}^{q-1}\omega^{-pj}
       \sum_{h=0}^{\ell}\binom{\ell}{h}(\log q)^{\ell-h}
       \left.\partial_s^h\Li_s(\omega^j)\right|_{s=0}.
 \label{agrc:Li-order-jets}
\end{align}
Together with the standard polygamma--Hurwitz relation
$\psi^{(m)}(x)=(-1)^{m+1}m!\zeta(m+1,x)$ for $m\geq1$,
this converts Theorem~\ref{agrc:beta-stieltjes} at positive rational
shifts into exact identities among harmonic-polygamma sums, generalized
Stieltjes constants, and order derivatives of finitely many polylogarithms.
Integer translation handles rational shifts outside $(0,1]$.

The Mellin kernel itself also has a multiple-polylogarithm interpretation:
\begin{equation}\label{agrc:Li-all-ones}
 \sum_{n\geq0}E_{n,r}z^n
  =\frac{[-\log(1-z)]^r}{r!(1-z)}
  =\frac{\dd}{\dd z}\Li_{\{1\}^{r+1}}(z),\qquad |z|<1.
\end{equation}
Thus the harmonic identities arise from Mellin transforms of actual
polylogarithmic derivatives, rather than from an analogy of notation.
Equation~\eqref{agrc:Hurwitz-Li} evaluates a prescribed finite average;
it does not imply reduction of its individual higher-depth analogues.

\section{Ordinary antiderivatives and Gamma reflection integrals}
\label{agrc:primitives}

\subsection{A log-Gamma and Hurwitz primitive}

Let $S_1(x)$ denote the ordinary sum in \eqref{agrc:first-log}, for real
$x>0$. An explicit antiderivative is
\begin{equation}\label{agrc:S1-primitive}
 \boxed{\displaystyle
 \int S_1(x)\dd x=
 \left(x-\frac12\right)\log\Gamma(x)-2\zeta'(-1,x)
 +\frac{x^2-x}{2}-\frac{x}{2}\log(2\pi)+\mathrm{constant}.}
\end{equation}
Indeed, differentiating the Hurwitz shift identity in the spectral
variable gives
\[
 \partial_x\zeta'(-1,x)=-\zeta(0,x)+\zeta'(0,x)
 =x-\tfrac12+\log\Gamma(x)-\tfrac12\log(2\pi).
\]
Direct differentiation of \eqref{agrc:S1-primitive} therefore gives
\eqref{agrc:first-log}. Normal convergence of the grouped series locally
in $x>0$ also permits ordinary integration term by term over a compact
interval. The formula is not a finite-part integral.

\subsection{Moving-centered reflection primitives at every depth}

For $c_u=(1+u)/2$, let
\[
 \Ccal_\pm(s,u;x)=\Bcal(s,u;c_u+x)\pm\Bcal(s,u;c_u-x).
\]
In the initial convergence domain, differentiation gives
\begin{equation}\label{agrc:parity-derivatives}
 \partial_x\Ccal_\pm(s,u;x)=-s\Ccal_\mp(s+1,u;x).
\end{equation}
It extends meromorphically and, for $s\ne1$, supplies the primitive
$-\Ccal_\mp(s-1,u;x)/(s-1)$. At $s=1$ this expression must be continued
with its pole terms; there is no license to substitute into an
indeterminate quotient.

At $s=2$ the anchored primitives are ordinary integrals. For real
$|x|<1/2$ and $u$ sufficiently small,
\begin{align}
 \int_0^x\Ccal_+(2,u;v)\dd v
   &=B(c_u-x,-u)-B(c_u+x,-u),\label{agrc:primitive-plus}\\
 \int_0^x\Ccal_-(2,u;v)\dd v
   &=2B(c_u,-u)-B(c_u+x,-u)-B(c_u-x,-u).
 \label{agrc:primitive-minus}
\end{align}
The apparent beta poles at $u=0$ cancel. These identities follow by
integrating $\Bcal(2,u;a)=-\partial_a B(a,-u)$ first for $\Re u<0$;
normal convergence in a neighbourhood of $u=0$ proves their extension.
All depth coefficients give finite Gamma/polygamma primitives.
For example, the depth-zero plus identity reduces to
\begin{equation}\label{agrc:trigamma-primitive}
 \int_0^x\bigl[\psi'(1/2+v)+\psi'(1/2-v)\bigr]\dd v
       =\pi\tan(\pi x).
\end{equation}

There is also a useful fixed-$a$ presentation of the moving endpoint:
\begin{align}
 \int_a^{1+u-a}\Bcal(2,u;x)\dd x
 &=B(a,-u)-B(1+u-a,-u)\notag\\
 &=\frac{\Gamma(a)\Gamma(-u)}{\Gamma(a-u)}
       \left(1-\frac{\sin\pi a}{\sin\pi(a-u)}\right).
 \label{agrc:moving-endpoint}
\end{align}
The ratio of the beta terms follows from Gamma reflection. For real
$0<a<1$ and small real $u$, the oriented integral is ordinary; the
complex-parameter statement is its analytic germ. At $u=0$ the right
side has a removable value $\pi\cot(\pi a)$. Higher coefficients include
the moving-endpoint contributions automatically.

\section{Several unequal frequencies and independent harmonic depths}
\label{agrc:multifrequency}

The parity mechanism is not restricted to one summation variable or a
single harmonic-depth parameter. Fix $d\geq1$ and positive real
frequencies $\omega_1,\ldots,\omega_d$. Put
\[
 \alpha_i=1+u_i,\qquad U=\sum_i u_i,
 \qquad a(\bm u)=\frac12\sum_i\alpha_i\omega_i.
\]
Define, initially in an absolute-convergence domain,
\begin{equation}\label{agrc:multi-B}
 \Bcal_{\bm\omega}(s,\bm u)=
 \sum_{n_1,\ldots,n_d\geq0}
 \frac{\prod_{i=1}^d(1+u_i)_{n_i}/n_i!}
      {\left(a(\bm u)+\sum_i\omega_i n_i\right)^s}.
\end{equation}
For instance, $\Re s>d+\sum_i|u_i|$ and all $u_i$ sufficiently small
provide a common domain. The denominator then has positive real part.

\begin{theorem}[Unequal-frequency all-depth reflection zeros]
\label{agrc:multi-zero}
Let $m\geq0$ be an integer with $d+m$ odd. The family
\eqref{agrc:multi-B} is jointly holomorphic near $(s,\bm u)=(-m,\bm0)$,
and
\begin{equation}\label{agrc:multi-zero-eq}
 \Bcal_{\bm\omega}(-m,\bm u)=0.
\end{equation}
Therefore every multi-depth coefficient $[\bm u^{\bm r}]$ has a
holomorphic zero there. No rational dependence among the frequencies
is required.
\end{theorem}
\begin{proof}
The product of binomial series yields
\begin{equation}\label{agrc:multi-Mellin}
 \Gamma(s)\Bcal_{\bm\omega}(s,\bm u)
 =\int_0^\infty t^{s-1}
             \prod_{i=1}^d\bigl(2\sinh(\omega_i t/2)\bigr)^{-1-u_i}
                                                      \dd t.
\end{equation}
Near zero, the kernel is
\[
 t^{-d-U}\prod_i\omega_i^{-1-u_i}
  \prod_i\left(\frac{\omega_i t}{2\sinh(\omega_i t/2)}\right)^{1+u_i}.
\]
The last product is an even analytic function of $t$, jointly analytic
in $\bm u$. Uniform Mellin subtraction therefore has possible moving
poles only at $s=d+U-2j$. At $s=-m,\bm u=0$, a pole would require
$2j=d+m$, excluded by the parity assumption. Choosing a sufficiently
large subtraction order proves joint holomorphy of the Mellin side
near the specified point. Multiplication by $1/\Gamma(s)$ then gives
the zero. Uniform local holomorphy permits every multi-depth extraction.
\end{proof}

For $d=1$ the protected points are $0,-2,-4,\ldots$; for $d=2$ they
are $-1,-3,-5,\ldots$. This is a statement about the family and all its
order derivatives, not just about a fixed integer-order Barnes zeta.

To display the harmonic content explicitly, let
$a_0=\frac12\sum_i\omega_i$ and, for a multi-index $\bm r$, define
\[
 Z_{\bm r}^{\bm\omega}(s)=
 \sum_{\bm n\geq0}\frac{\prod_i E_{n_i,r_i}}
                         {(a_0+\sum_i\omega_i n_i)^s}.
\]
Taylor expansion of the moving center gives the finite exact formula
\begin{equation}\label{agrc:multi-coeff}
 [\bm u^{\bm r}]\Bcal_{\bm\omega}(s,\bm u)
 =\sum_{\bm k\leq\bm r}(-1)^{|\bm k|}(s)_{|\bm k|}
       \prod_i\frac{(\omega_i/2)^{k_i}}{k_i!}
          Z_{\bm r-\bm k}^{\bm\omega}(s+|\bm k|).
\end{equation}
Theorem~\ref{agrc:multi-zero} gives all-order cancellation identities
for these combinations. At a singular specialization, the complete
combination is continued before its value is taken.

If the center is instead displaced by a fixed $x$, reflection
$x\mapsto-x$ again separates the even and odd heat coefficients.
The plus combination has the parity class of Theorem~\ref{agrc:multi-zero};
the minus combination has the opposite class of nonpositive integers.
This follows by multiplying the even heat factor by $2\cosh(xt)$ or
$-2\sinh(xt)$, respectively.

The result does not say that all unequal-frequency coefficients reduce
to ordinary zeta constants. For noncommensurate frequencies, the
underlying Barnes-type periods remain distinct coordinates. The exact
zero statement is independent of such a reduction problem.

\section{Audit, necessary corrections to extensions, and verification}
\label{agrc:audit-section}

\subsection{The fixed-shift negative Laurent theorem is prior work}

The preserved source \cite{repo-laurent} already proves the following
local rule (its outer-shift notation differs by one). With our heat
coefficients and $m\geq0$, put
\[
 \frac{c_{m+1}(u,a)}{\Gamma(u-m)}
       =\sum_{j\geq1}\lambda_{m,j}(a)u^j.
\]
Then
\begin{equation}\label{agrc:prior-Laurent}
 Z_r(-m+\varepsilon,a)
 =\sum_{h=0}^r\lambda_{m,r+1-h}(a)\varepsilon^{-h}
        +O(\varepsilon).
\end{equation}
A short explanation is worth retaining: near $s=-m$, only the heat
coefficient $c_{m+1}$ has a moving denominator $\varepsilon-u$.
All other terms are holomorphic and, after multiplication by
$1/\Gamma(-m+\varepsilon)$, are $O(\varepsilon)$ at every fixed depth.
Multiplying the remaining numerator by the reciprocal Gamma series
produces exactly \eqref{agrc:prior-Laurent}.

This is not a new result of the present continuation. It is included
to normalize examples and to make the additional moving-reflection
condition transparent. The analytic remainder starts at
$O(\varepsilon)$, so it can and generally does affect positive spectral
 derivatives. In particular, at depth zero these include
$\zeta'(-m,a)$, and at $m=0$ they already include $\log\Gamma(a)$.

\subsection{A version-specific indexing correction in a cited preprint}

The unnumbered Stieltjes-limit display at the bottom of page 1 of
\cite{Can2024}, specifically arXiv version 1, starts its sum at $n=1$
while using the Hurwitz normalization with $n=0$. The PDF page was
checked directly. For real $a>0$, the corrected formula is
\begin{equation}\label{agrc:corrected-Stieltjes-limit}
 \gamma_m(a)=\lim_{N\to\infty}\left(
 \sum_{n=0}^{N}\frac{\log^m(n+a)}{n+a}
                  -\frac{\log^{m+1}(N+a)}{m+1}\right).
\end{equation}
If the lower index $n=1$ is retained, the missing term
$\log^m(a)/a$ must be added. Indeed,
$\zeta(s,a+1)=\zeta(s,a)-a^{-s}$ gives
\[
 \gamma_m(a+1)=\gamma_m(a)-\frac{\log^m a}{a}.
\]
The printed $n=1$ limit is $\gamma_m(a+1)$, not $\gamma_m(a)$.
Already at $m=0,a=1$, it becomes
$\lim_{N\to\infty}(H_{N+1}-1-\log(N+1))=\gamma-1$.
This identifies a specific off-by-one error; it is not an objection to
that paper's other results, and no claim is made about later revisions.
The convention \eqref{agrc:Stieltjes} and all identities in this report
use the corrected normalization.

\subsection{Three invalid extensions, with exact corrections}

No contradiction to the cited Gauss subtraction or fixed-shift Laurent
theorems was found in the targeted audit. The following are explicit
corrections to tempting extensions, not allegations that those sources
made the errors.

\paragraph{Frozen reflection is not moving reflection.}
Replacing $1+u-a$ by $1-a$ leaves the first heat coefficient sum equal
to $u$. As a concrete consequence,
\begin{equation}\label{agrc:frozen-counterexample}
 Z_1(s,a)+Z_1(s,1-a)=1+O(s)\qquad(s\to0),
\end{equation}
for $0<\Re a<1$. The sum is holomorphic but is not zero. Its poles
cancel because their leading heat coefficients are opposite; its
constant terms sum to one. The extra shift-derivative term in
\eqref{agrc:Rcoeff} is exactly what removes this defect.

\paragraph{Normal jets are not determined by a resonant restriction.}
For $a=b=1$ and $s=-1+\varepsilon$, the exact polynomial is
\begin{equation}\label{agrc:normal-counterexample}
 A_1(1,1,-1+\varepsilon;c)
          =(c-\tfrac12)\varepsilon-\tfrac12\varepsilon^2.
\end{equation}
Although its resonant value is zero and independent of $c$, its first
normal derivative is $c-1/2$. Formula~\eqref{agrc:Gauss-jets} therefore
needs actual normal derivatives; the restriction
\eqref{agrc:residue-formula} cannot replace them.

\paragraph{Finite-part extraction does not commute with multiplication
by a varying coefficient without residue terms.}
For a Laurent series $V(\varepsilon)$,
$\FP(\varepsilon V)=\Res V$, not zero in general.
Equation~\eqref{agrc:Z2-fp} displays the resulting correction explicitly.
Similarly, omitting the second subtraction in \eqref{agrc:Dx} destroys
the normal convergence needed for its all-order derivative formulas.

\subsection{Reproducible checks and their limits}

The supplied verifier uses two independent coefficient constructions:
Bernoulli recursions for $A_j$ and finite Stirling sums for $Q_{r,k}$.
The exact run contains 436 identities: 16 fully symbolic resonant-tail
checks, 189 higher-order rational checks, 5 comparison-shift identities,
25 Appell identities, 154 finite-product collapse checks, 13 moving
reflection checks, 10 centered parity checks, 15 unequal-frequency
parity checks, and 9 Gamma-cumulant checks. Four negative controls reject
specified corruptions of a tail sign, a reflection shift, a center slope,
and a Stieltjes sign.

The numerical run uses 65 decimal digits and makes 30 comparisons.
These include ordinary Mellin integrals versus the Gamma quotient,
log-sinh moments versus the independent cumulant recursion,
subtracted Gamma-ratio sums versus \eqref{agrc:Dx-closed}, direct
Hurwitz differentiation of \eqref{agrc:S1-primitive}, and Cauchy
coefficient extraction for the second Stieltjes jet. The largest
reported normalized discrepancy is approximately $1.99\times10^{-51}$.
The Gamma-ratio series accelerator uses finite asymptotic tails; these
floating-point diagnostics are \emph{not} rigorous interval enclosures.

The finite tests do not prove the arbitrary-index theorems. Those are
proved in the preceding sections. Conversely, a passing finite test
should not promote a different pending source claim to proved status.
There is no proof-assistant formalization in this package.

\section{Further research questions and concrete next targets}
\label{agrc:future}

\subsection{Normal spectral jets of the resonant tail}
Theorem~\ref{agrc:tail} gives a closed finite formula for the full tail
on $s=-r$. Find a similarly explicit representation of
$\partial_s^\ell A_{r+k}|_{s=-r}$ for arbitrary $\ell$. A useful first
result would be a hypergeometric inhomogeneous differential equation
for the Borel transforms at $\ell=1,2$, together with normalized initial
data. Equation~\eqref{agrc:normal-transport} is a necessary compatibility
condition. A solution would make the Stieltjes identities
\eqref{agrc:Gauss-jets} substantially more explicit beyond their
constant coefficient.

\subsection{Derivatives at the protected common zeros}
The common zero $C_r(-2m)=0$ is local and exact. Its derivative generally
retains the analytic Mellin remainder. Determine which finite combinations
of $C_r'(-2m)$ eliminate that remainder, and identify the natural
coordinates for the first surviving cases, beginning with $C_1'(0)$
and $C_1'(-2)$. The target is an additional identity, not a numerical
claim that a selected constant basis is complete.

\subsection{Colored moving reflection and the Gaussian targets}
Replace the kernel by products involving $1-\xi\e^{-t}$ for roots of
unity $\xi$. Derive the exact reflected colors, branches, and local
normalization before extracting harmonic derivatives. A concrete goal
is a functional identity whose Gaussian specialization supplies relations
outside the already insufficient restricted S6 presentation. Neither a
small numerical residual nor the present uncolored centering theorem
settles S6 or S8. Their exact frozen coefficient vectors must be retained
when testing any new relation.

\subsection{Commensurate versus noncommensurate frequencies}
For positive rationally related $\omega_i$, determine how much of
\eqref{agrc:multi-coeff} reduces by finite character filters and
partial fractions. Compare this with genuinely irrational frequency
ratios. The common-zero theorem is uniform in that distinction, but
closed evaluations of its derivatives need not be. A first concrete
case is two frequencies in ratio $1:2$ at multi-depth $(1,1)$.

\subsection{All-order reduction of the shifted logarithmic family}
Theorem~\ref{agrc:beta-stieltjes} expresses every logarithmic order in
Gamma/polygamma, generalized Stieltjes, and Hurwitz spectral coordinates.
At rational shifts, determine additional identities among the order
 derivatives of polylogarithms in \eqref{agrc:Li-order-jets}. Such a
reduction should be specified in a fixed finite coordinate space and
proved; it should not be inferred from an unsuccessful relation search.

\subsection{Connection identities for the Borel hypergeometric germ}
The function $\e^{ct}{}_2F_1(a,b;r+1;1-\e^t)$ supplies an explicit analytic
germ for every resonant tail. Determine useful connection identities
between this germ, the original Gamma ratio, and logarithmic expansions
of Gauss's function at its singular endpoint. The objective is exact
functional continuation and coefficient transport, not merely estimates
for optimal truncation.

\subsection{Formalization and integration}
The finite proof of Theorem~\ref{agrc:residue} is a natural first formal
 target: it needs Bernoulli finite differences, polynomial divisibility,
a degree bound, and one leading coefficient. The tail theorem then adds
formal Laurent transforms and uniqueness for a power-series differential
equation. For the analytic side, reusable lemmas for normally convergent
parameter families, Mellin subtraction, and meromorphic coefficient
extraction would support all-depth statements without separate proofs
for every tabulated case.

\appendix
\section{Algorithms and normalization checks}
\label{agrc:algorithms}

\subsection{Finite construction of every resonant coefficient}

For fixed integers $r,k\geq0$, \eqref{agrc:tail-formula} can be evaluated
without generating any $A_j$ of index below $r$:
first form $P_r(a)P_r(b)$; then compute the Stirling triangle through row
$k$; finally evaluate the two finite sums in \eqref{agrc:Q-definition}.
The construction uses only rational arithmetic when $a,b,c$ are rational.
At a root of either Pochhammer factor, use the polynomial formula
 directly, not division by the factor. Differentiation at those roots
uses \eqref{agrc:root-jets} or direct polynomial differentiation.

For an independent implementation, construct $D_1,\ldots,D_{r+k}$ from
Bernoulli polynomials and use \eqref{agrc:A-rec}. Equality of these two
independent constructions is the main exact regression test.
The formulas for $Q_{r,1}$ and $Q_{r,2}$ in
\eqref{agrc:Q-examples} are useful low-order normalization checks.

\subsection{Gamma coefficient recursions}

If $F(u)=\exp(\sum_{j\geq1}\eta_j u^j/j)=\sum f_ru^r$, then
\begin{equation}\label{agrc:general-coeff-rec}
 f_0=1,\qquad rf_r=\sum_{j=1}^r\eta_j f_{r-j}.
\end{equation}
For the centered Gamma quotient $F_c$, use
\[
 \eta_1=-2L,\qquad
 \eta_j=\begin{cases}
 \zeta(j),&j\text{ even},\\
 (2^{1-j}-1)\zeta(j),&j\geq3\text{ odd}.
 \end{cases}
\]
For $G(u)=1/\Gamma(1+u)$, use
$\eta_1=\gamma$ and $\eta_j=(-1)^{j+1}\zeta(j)$ for $j\geq2$.
For $C_r(2)/(\pi^2/2)$, use \eqref{agrc:kappa}.
These three recursions must not be conflated: their even zeta
coefficients differ.

\subsection{Notation crosswalk for repository integration}

Our $Z_r(s,a)$ is the family denoted $E_r(s;a-1)$ in
\cite{repo-laurent}. Its heat coefficient $b_j(u,a-1)$ is our
$c_j(u,a)$. Our $A_j(a,b,s;c)$ agrees with the Gauss coefficient in
\cite{repo-gauss}, and is unrelated to that source's generalized harmonic
functions with a similar letter. The centered coefficient $C_r(s)$ is
new notation local to this report; it must not be identified with the
fixed-shift Gamma polynomial sometimes denoted $C_r(a)$ elsewhere in
the manuscript. All theorem and equation labels in this article use
the prefix \texttt{agrc:}.

\section{Source and artifact boundaries}

The source audit pins the relevant repository content by commit and blob
identifiers in the file \path{audit/source_manifest.json}. It records the exact
research question addressed and distinguishes the prior negative-Laurent
 theorem from the moving-reflection extension. Only inspected textual
sources are used as repository evidence; there is no claim to have
replayed every incoming ZIP.

The archive contains the standalone article, its shared TeX body, a
verification script, exact tables, exact and numerical run records,
an integration guide, and a build script. The article can be rebuilt
with a standard LaTeX installation; the verification suite uses SymPy
and mpmath. Generated log and auxiliary files are not mathematical
certificates. A checksum manifest records the delivered artifacts.
